\chapter{希尔伯特零点定理}

定理 1. --- 设 A 是一个整环，K 是它的分式域，L 是一个交换 K-代数。假设 A-代数 L 由有限个元素生成，且 L 是一个体。

\begin{enumerate}
\def\labelenumi{\alph{enumi})}
\item
  L 在 K 上的次数有限。
\item
  存在一个非零元素 \(a\)（属于 \(\mathrm{A}\)），使得 \(\mathrm{K}\) 等于 \(\mathrm{A}[a^{-1}]\)。
\end{enumerate}

设 \(S\) 是 A-代数 \(L\) 的一个生成子集。我们对 \(S\) 的基数用归纳法。

首先，假设 S 的元素在 K 上都是代数的。这时 L 在 K 上的次数有限（V, §3, No.~2, 第 18 页，定理 2）。设 \((e_i)_{i\in I}\) 是 L 在 K 上的一个基。存在一个非零元素 \(a\)（属于 A），使得 S 的元素、元素 1 以及元素 \(e_i e_j\)（\(i\)、\(j\) 属于 I）关于这个基的坐标都属于 \(\mathrm{A}[a^{-1}]\)。于是诸线性组合 \(\sum_{i\in I}a_i e_i\)（其中 \(a_i\in \mathrm{A}[a^{-1}]\)，对每个 \(i\)）所成的集合是 L 的一个子环。按构造它包含 A 与 S，从而等于 L。特别地，它包含 \(\mathrm{Ke}_1\)，故 K 等于 \(\mathrm{A}[a^{-1}]\)。

现在设一个元素 \(s\)（属于 \(S\)）在 \(K\) 上是超越的。用 \(E\) 表示环 \(A[s]\) 的分式域。\(A[s]\)-代数 \(L\) 由 \(S - \{s\}\) 生成。由归纳假设，存在一个非零多项式 \(P \in A[X]\)，使得 \(E\) 等于 \(A[s][P(s)^{-1}]\)。设 \(\overline{K}\) 是 \(K\) 的一个代数闭包（V, §4, No.~3, 第 23 页，定理 2）。由于体 \(\overline{K}\) 是无限的（V, §4, No.~1, 第 20 页，命题 3），存在一个元素 \(x\)（属于 \(\overline{K}\)），使得 \(P(x) \neq 0\)。设 \(\varphi : E \to \overline{K}\) 是从 K-代数 \(E = A[s][P(s)^{-1}]\) 到 K-代数 \(\overline{K}\) 的把 \(s\) 映为 \(x\) 的唯一同态。这是荒谬的，因为 \(E\) 是 \(K\) 的一个超越扩张，而 \(\overline{K}\) 是 \(K\) 的一个代数扩张。这就完成了定理的证明。

推论 1. --- 设 K 是一个交换域，A 是一个由有限个元素生成的交换 K-代数，m 是 A 的一个极大理想。a) A/m 在 K 上的次数有限。

\begin{enumerate}
\def\labelenumi{\alph{enumi})}
\setcounter{enumi}{1}
\tightlist
\item
  设 \(\Omega\) 是 K 的一个代数闭扩张。存在一个从 A 到 \(\Omega\) 的以 \(\mathfrak{m}\) 为核的 K-代数同态。
\end{enumerate}

K-代数 A/m 由有限个元素生成，且 A/m 是一个体。由定理 1，A/m 的次数有限；由此得断言 a)。K 的每个次数有限的扩张都同构于 \(\Omega\) 的一个子扩张（V, §4, No.~1, 第 20 页，定理 1）；这就给出 b)。

推论 2. --- 设 K 是一个交换域，n 是一个自然数，\((\mathrm{P}_{i})_{i\in\mathrm{I}}\) 是 \(K[X_{1},\ldots,X_{n}]\) 的元素的一个族，\(\Omega\) 是 K 的一个代数闭扩张。下列条件等价：

\begin{enumerate}
\def\labelenumi{(\roman{enumi})}
\item
  多项式 \(\mathrm{P}_i\) 在 \(\Omega^n\) 中没有公共零点。
\item
  存在一个有限支撑族 \((\mathrm{Q}_i)_{i\in \mathrm{I}}\)（由 \(\mathrm{K}[\mathrm{X}_1,\dots ,\mathrm{X}_n]\) 的元素组成），使得 \(\sum_{i\in \mathrm{I}}\mathrm{P}_i\mathrm{Q}_i = 1\)。
\end{enumerate}

用 A 表示环 \(K[X_{1},\ldots,X_{n}]\)，用 a 表示由多项式 \(P_{i}\) 生成的理想。条件 (i) 表示不存在从 A/a 到 \(\Omega\) 的 K-代数同态。若它被满足，则由推论 1，环 A/a 没有极大理想，从而是零环，且 1 属于 a。这证明了 (i) 蕴含 (ii)。(ii) \(\Rightarrow\) (i) 是明显的。

